A packing of $\PG(3,q)$ 
is a set of pairwise line-disjoint spreads of $\PG(3,q)$ of size $q^2+q+1.$ 
Each spread contains $q^2+1$ lines. 
A simple counting argument shows that 
every line is contained in exactly one spread of the packing.
The classification problem for packings 
is the problem of determining a complete set of pairwise non-isomorphic 
packings. 
Orbiter can be used to classify packings for small parameters. 
It is sometimes useful to make a symmetry assumption.
This means that only those packings will be found that satisfy the symmetry assumption. 
The reason for making such an assumption is that the problem becomes easier and hence more tractable.
Often, an assumption is made that the packings are invariant under a (nontrivial) group $H$.
This section describes various ways in which Orbiter can help find and classify packings, 
with or without symmetry assumption.

\bigskip


Table~\ref{tab:packings:was} list Orbiter commands related 
to the construction of packings with assumed symmetry.
\orbitertableofcommands{packing_was}


\bigskip

Table~\ref{tab:packing:classify} list Orbiter activities for the classification of packings.
\orbitertableofcommands{packing_classify}



\bigskip

Table~\ref{tab:packings:was:activity} list Orbiter activities for packings with assumed symmetry.
\orbitertableofcommands{packing_was_activity}


\bigskip

Table~\ref{tab:packings:was:fixpoints:activity} list Orbiter activities for packings with assumed symmetry and with fixpoints.
\orbitertableofcommands{packing_was_fixpoints_activity}

\bigskip

Table~\ref{tab:packings:long:orbits} list Orbiter commands related to the construction 
of packings with assumed symmetry by picking long orbits.
\orbitertableofcommands{packing_long_orbits_description}


\bigskip



There are two isomorphism types of spreads in $\PG(3,3)$. 
The following command computes the orbits of each and 
creates a table of all (labeled) spreads in $\PG(3,3).$ 


\sourcecode{PG_3_3_spread_table}


\bigskip

There are 21 isomorphism types of spreads in $\PG(3,5).$ 
The regular spread has Orbiter catalogue number equal to 12. 
The following command creates a table of all regular spreads in $\PG(3,5)$:

\sourcecode{spread_table_PG_3_5_regular}

There are $155000$ regular spreads.

\bigskip

Suppose we are interested in spreads of $\PG(3,3)$ invariant under a group of order $5.$ 
To do so, we can apply a sequence of Orbiter commands as follows.
To begin with, we need to find a generator for a subgroup of order 5 in $\PGL(4,3).$
To this end, we compute the conjugacy classes of the group, using the following command:


\sourcecode{PGL_4_3_classes}

The actual computation is done in Magma, 
but Orbiter provides us the group as matrix group and produces a latex report.
The report contains all the information we need. 
At first, we pick the representative of the unique class of order $5$ 
(note that Magma does not behave deterministically, 
so different runs of Magma may result in different elements of order 5):

\sourcecodevariable{PGL_4_3_CLASS_5A_REP}

Let $H$ be the subgroup generated by this element. 
We also need the generators for the normalizer $N$ of $H$. 
The subgroup $N$ has order 160, 
and generators can be found in the report on conjugacy classes as well. 
Using copy-paste, we store the generators in a makefile variable like so:

\sourcecodevariable{PGL_4_3_CLASS_5A_N_GENS}


The next step is to prepare for a search for all packings 
invariant under the chosen group $H$ of order 5.
To this end, we apply the following command:

\sourcecode{PG_3_3_assume_5_graph}


It produces the orbits of the group $H$ on packings, filtering out those orbits 
which cannot be used for packings because the spreads are not line disjoint. 
It then classifies the orbits by length. The only possible lengths are one and 5. 
The orbits of length one are called the fixpoints. 
Because we need 13 spreads in a packing, we must have at least 3 fixed points 
(recall that point means spread).
The command classifies the ways in which one can pick 3 
fixed points up to the action of $N$. 
It turns out that there are three ways, up to isomorphism. 
This creates three disjoint cases.
For each case, we need to continue the construction by adding two orbits of length 5.
In order to do so, we consider each of the three cases separately. 
In each case, we first filter the set of orbits of length 5 according 
to wether they are compatible with the fixed sread configuration. 
This means that the spreads in the orbit are supposed 
to be line disjoint from the spreads in the three fixed spreads. 
Once done, we compute a graph. 
The vertices of the graph are the orbits of spreads of length 5 
that are compatible with the fixed point configuration. 
Two vertices are adjacent in the graph 
if the associated spread orbits 
are compatible, which means that they are line-disjoint. 
The packings that contain the fixed configuration and 
consist of an additional set of 2 orbits of length 5 
are in one-to-one correspondence to the cliques of size 2 in the graph.
We employ Orbiter's clique finding algorithm (see Section~\ref{sec:graphtheory:cliques}) 
to compute all these cliques (in each of the three cases):

\sourcecode{PG_3_3_assume_5_fpc0_lo_cliques}

\sourcecode{PG_3_3_assume_5_fpc1_lo_cliques}

\sourcecode{PG_3_3_assume_5_fpc2_lo_cliques}

Once we are done with the cliques, 
we need to assemble the packings that arise 
by joining the orbits of length 5 associated 
with the elements of the clique to the fixed point configuration. 
Once this is done, we assemble the packings:

\sourcecode{PG_3_3_assume_5_read_solution}


As it turns out, the clique finding commands find $0+4+10=14$ cliques, 
for a total of 14 packings. However some of these packings can be isomorphic. 
For this reason, we  perform a final isomorphism test, based on Nauty. 
We use Nauty through the Orbiter interface 
(for more on that, see Chapter~\ref{chapter:canonical:forms}). 
The following command does the job:



\sourcecode{PG_3_3_assume_5_packings_classify}



We end up with a list of isomorphism types of packings in $\PG(3,3)$ 
invariant under the chosen group $H$ of order 5.
Finally, the command prepares a latex report for these packings.






\bigskip

%Suppose we are interested in packings invariant under a group of order 31. 
%We pick the matrix
%$$
%\left[
%\begin{array}{*{4}{r}}
%1 & 0 & 0 & 0\\
%0 & 3 & 4 & 3\\
%0 & 3 & 3 & 4\\
%0 & 3 & 2 & 3\\
%\end{array}
%\right]
%$$
%of order 124 in $\GL(4,5).$
%At first, we store the element in a makefile variable:
%\sourcecode{PGL_4_5_SUBGROUP_31_ME}

%We compute the normalizer of the cyclic subgroup generated by this element inside $\GL(4,5)$:
%\sourcecode{PG_3_5_element_of_order_31_GL_normalizer}

%We compute the normalizer of the cyclic subgroup generated by this element inside $\PGL(4,5)$:
%\sourcecode{PG_3_5_element_of_order_31_ME_normalizer}
%We store the normalizer in a makefile variable:
%\sourcecode{PGL_4_5_SUBGROUP_31_ME_NORMALIZER}

%We create packings invariant under the cyclic group of order 31 generated by this element. 
%\sourcecode{PG_3_5_assume_31_graph}




